% Why the walk cannot be bounded by the querying box's own footprint: two
% overlapping boxes whose minimum-corner cells are incomparable.
%
% Conventions as in tikz/cellmin.tex -- no prose inside the drawing, grids
% drawn line by line, and the geometry is real: cells are 0.62 wide, five by
% three, numbered row-major from the bottom left, and every cell index below
% follows from the coordinates in \boxA and \boxB.
\begin{tikzpicture}[
  scale=1.65,
  font=\scriptsize,
  lbl/.style={font=\scriptsize, text=sccdInk, inner sep=1.5pt},
  tiny/.style={font=\tiny, text=sccdInk, inner sep=1pt},
  ttl/.style={font=\scriptsize\bfseries, text=sccdInk, inner sep=0pt},
  idle/.style={draw=sccdInk!22, line width=0.5pt},
  qbox/.style={draw=sccdTeal, line width=1.0pt},
  pbox/.style={draw=sccdPlum, line width=1.0pt},
  read/.style={fill=sccdTeal!14},
  unread/.style={fill=sccdAmber!22},
]

% A: minimum corner in column 2, row 0.  B: column 1, row 1.
\def\boxA{1.50/0.10/1.90/0.90}
\def\boxB{0.80/0.70/2.20/1.00}

\newcommand{\cellfill}[3]{%
  \fill[#3] ({#1*0.62},{#2*0.62}) rectangle ({(#1+1)*0.62},{(#2+1)*0.62});}

\newcommand{\scenegrid}{%
  \foreach \x in {0,0.62,1.24,1.86,2.48,3.10}
    \draw[draw=sccdInk!30, line width=0.3pt] (\x,0) -- (\x,1.86);
  \foreach \y in {0,0.62,1.24,1.86}
    \draw[draw=sccdInk!30, line width=0.3pt] (0,\y) -- (3.10,\y);
  \draw[draw=sccdInk!45, line width=0.4pt] (0,0) rectangle (3.10,1.86);}

\newcommand{\cellnumbers}{%
  \foreach \r in {0,1,2}
    \foreach \c in {0,1,2,3,4} {
      \pgfmathtruncatemacro{\id}{\r*5+\c}
      \node[tiny, text=sccdInk!45, anchor=south west, fill=white,
            inner sep=0.5pt] at ({\c*0.62+0.03},{\r*0.62+0.03}) {\id};}}

% The two boxes, with the minimum corner each is binned at.
\newcommand{\thetwoboxes}{%
  \foreach \xa/\ya/\xb/\yb in \boxA {
    \draw[qbox] (\xa,\ya) rectangle (\xb,\yb);
    \fill[sccdTeal] (\xa,\ya) circle (1.4pt);}
  \foreach \xa/\ya/\xb/\yb in \boxB {
    \draw[pbox] (\xa,\ya) rectangle (\xb,\yb);
    \fill[sccdPlum] (\xa,\ya) circle (1.4pt);}}

% ========================= (a) the two boxes =================================
\begin{scope}
  \node[ttl, anchor=south west] at (0,2.02) {(a) $A$ and $B$ overlap};

  \scenegrid
  % The overlap itself, so that the reader can see there is a pair to find.
  \fill[sccdInk!12] (1.50,0.70) rectangle (1.90,0.90);
  \thetwoboxes
  \cellnumbers
  \node[lbl, text=sccdTeal, anchor=north west] at (1.92,0.45) {$A$};
  \node[lbl, text=sccdPlum, anchor=south west] at (2.22,1.00) {$B$};
\end{scope}

% ===================== (b) A reads its own footprint =========================
% A is binned in cell 2 and its maximum corner in cell 8, so its footprint is
% columns 2-3 crossed with rows 0-1: cells 2, 3, 7, 8. B is binned in cell 6,
% one column to the left of that, and is never read.
\begin{scope}[shift={(3.72,0)}]
  \node[ttl, anchor=south west] at (0,2.02) {(b) $A$'s footprint misses $B$};

  \cellfill{2}{0}{read}  \cellfill{3}{0}{read}
  \cellfill{2}{1}{read}  \cellfill{3}{1}{read}
  \cellfill{1}{1}{unread}
  \scenegrid
  \thetwoboxes
  \cellnumbers
  \node[lbl, text=sccdTeal, anchor=north west] at (1.92,0.45) {$A$};
  \node[lbl, text=sccdPlum, anchor=south west] at (2.22,1.00) {$B$};
\end{scope}

% ===================== (c) B reads its own footprint =========================
% B is binned in cell 6 and its maximum corner in cell 8, so its footprint is
% columns 1-3 on row 1: cells 6, 7, 8. A is binned in cell 2, one row below,
% and is never read.
\begin{scope}[shift={(0,-2.60)}]
  \node[ttl, anchor=south west] at (0,2.02) {(c) $B$'s footprint misses $A$};

  \cellfill{1}{1}{read}  \cellfill{2}{1}{read}  \cellfill{3}{1}{read}
  \cellfill{2}{0}{unread}
  \scenegrid
  \thetwoboxes
  \cellnumbers
  \node[lbl, text=sccdTeal, anchor=north west] at (1.92,0.45) {$A$};
  \node[lbl, text=sccdPlum, anchor=south west] at (2.22,1.00) {$B$};
\end{scope}

% ===================== (d) the walk that does find it ========================
% A holds the smaller linear index, so it reports the pair. Its own row starts
% at its own cell; the row above starts where the row's prefix maximum says
% something reaches back, which is column 1 -- left of A's own column.
\begin{scope}[shift={(3.72,-2.60)}]
  \node[ttl, anchor=south west] at (0,2.02) {(d) $A$ reads cells $2,3$ then $6,7,8$};

  \cellfill{2}{0}{read}  \cellfill{3}{0}{read}
  \cellfill{1}{1}{read}  \cellfill{2}{1}{read}  \cellfill{3}{1}{read}
  \scenegrid
  \thetwoboxes
  \cellnumbers
  \node[lbl, text=sccdTeal, anchor=north west] at (1.92,0.45) {$A$};
  \node[lbl, text=sccdPlum, anchor=south west] at (2.22,1.00) {$B$};
\end{scope}

% ================================== key ======================================
\begin{scope}[shift={(0,-3.30)}]
  \fill[sccdInk] (0.11,0.22) circle (1.4pt);
  \node[lbl, anchor=west] at (0.26,0.22) {the minimum corner a box is binned at};
  \fill[read] (0.00,-0.08) rectangle ++(0.22,0.16);
  \draw[draw=sccdInk!30, line width=0.3pt] (0.00,-0.08) rectangle ++(0.22,0.16);
  \node[lbl, anchor=west] at (0.26,0.00) {the cells the query reads};
  \fill[unread] (0.00,-0.38) rectangle ++(0.22,0.16);
  \draw[draw=sccdInk!30, line width=0.3pt] (0.00,-0.38) rectangle ++(0.22,0.16);
  \node[lbl, anchor=west] at (0.26,-0.30) {the partner's cell, left unread};
\end{scope}
\end{tikzpicture}
